\section{Fluid and rigid body}\label{sec_fl_rigid}

In this Section, we will discuss the interaction between a  \ac{bflu} and the simplest structure---a  \ac{rig} which is only allowed to turn about one point.


In all the systems that we will consider, we will introduce the \ac{ref} and \ac{cur} configuration \cite{kamenskyLectureNotesMAE,landauTheoryElasticity1986,slaughterLinearizedTheoryElasticity2002}, see \cref{fig_ref_curr_rigid}. All quantities relative to the \ac{ref} and \ac{cur} configuration will be denoted by the superscript \ac{ref} and \ac{cur}, respectively.
In the \ac{ref} configuration, the region occupied by the bulk fluid is described by Cartesian coordinates $\bxr$. The  axis of the ellipse, e..g, the rigid body, lies parallel to the $\xc^1$ axis.
In the \ac{cur} configuration, the region occupied by the bulk fluid is described by Cartesian coordinates $\bxc$. The  axis of the ellipse, e..g, the rigid body, forms an angle  $\theta$ (red arc) with respect to the $\xc^1$ axis.


The \ac{ref} configuration, is related to the \ac{cur} one as follows \cite{kamenskyLectureNotesMAE}, see \cref{fig_ref_curr_rigid}. At  time $t$, the point $\bxc$  corresponds to a point $\bxr$ in the reference configuration through the deformation field \cite{landauFluidMechanics1987}
\be\label{eq_def_u}
\budom^t(\bxr) \equiv \bxc - \bxr.
\ee
and the deformation-gradient tensor
\be
\label{eq_def_F}
F_{\alpha \beta}(\bu) \equiv \delta_{\alpha \beta} + \pderr{u^ \alpha}{\beta},
\ee
where we set
\begin{align}\label{eq_def_pder_ref}
        \partial^\reference_\alpha & \equiv \pder{}{\xr^\alpha}, \\
        \label{eq_def_pder_cur}
        \partial^\current_\alpha   & \equiv \pder{}{\xc^\alpha}.
\end{align}

Here, we will denote the mapping between $\bxc$ and $\bxr$ in \cref{eq_def_u} with the fields
\begin{align}\label{eq_def_phi}
        \bxc & = \bm{\varphi}^t(\bxr), \\
        \label{eq_def_psi}
        \bxr & = \bm{\psi}^t(\bxc).
\end{align}

We denote the \ac{bflu} velocity and surface tension by $\vfc(\bxc)$ and $\sigmafc(\bxc)$, respectively, where the superscript \ac{cur} specifies that these fields depend on the coordinate $\bxc$ in the \ac{cur} configuration.


The solution of the interaction dynamics between \ac{bflu} and \ac{rig} is  involved because of the presence of the moving boundary \pomellipsec, cf. \cref{fig_ref_curr_rigid}. Following the \ac{ale} procedure, we will  bridge between the Lagrangian and the Eulerian description used to describe, respectively,  \ac{rig} and \ac{bflu}  \cite{kamenskyLectureNotesMAE}. To achieve this, in what follows we will write the equations of motion for \ac{rig} and \ac{bflu}, which are naturally formulated the \ac{cur} configuration first, and then rewrite them in the \ac{ref} configuration.


\begin{figure*}
        \centering
        \includegraphics[width=\textwidth]{figures/figure_6/figure_6.pdf}
        \caption{
                \label{fig_ref_curr_rigid}
                \Acl{ref} and \acl{cur} configuration for the interaction between a bulk fluid and a rigid body.
                %
                \plab{A} \Acl{ref} configuration. The region $\omr$ is given by the rectangle $0 \leq \xr^1 \leq L$, $0 \leq \xr^2 \leq h$ with the elliptical hole, and is delimited by the mesh (gray lines).  Boundaries in the \acl{ref} configuration are denoted by $\pomineqr$, $\pomouteqr$, $\pomtopeqr$, $\pombottomeqr$ and $\pomellipseeqr$ (colored dashed lines). The ellipse focal point $\bfocalpoint$ (light green dot)  is also shown.
                %
                \plab{B} \Acl{cur} configuration. The region $\omc$ is delimited by the mesh (gray lines). Boundaries in the \acl{cur} configuration are denoted by $\pomineqc$, $\pomouteqc$, $\pomtopeqc$, $\pombottomeqc$ and $\pomellipseeqc$ (colored dashed lines), and the rotational angle $\theta$ of the rigid body is shown as an arc (black dotted line). The Cartesian coordinates $\bxc$ are also shown.
        }
\end{figure*}

%M revise to embed in text
Details on the motion of \ac{rig}, \ac{bflu} and \ac{dom} can be found in \cref{sec_rigid_body_body,sec_rigid_body_fl,sec_eq_fluid_domain}, respectively, and the temporal discretization is discussed in \cref{sec_temp_discretization}.



We can now iterate in time for the ensemble of fields as follows. At each step $n$, given $\theta^{n-1}$, $\omega^{n-1}$, ${\bm \vfr}^{n-1}$, $\sigmafr^{n-3/2}$, $\budom^{n-1}$ and $\budomdot^{n-1}$ from the preceding step, we
\begin{enumerate}
        \titleditem{Update \ac{rig}}{ \label{fl_rigid_step_1}
                Solve the algebraic equations \cref{eq_def_omega_disc,eq_fl_rigid_rigid_reference_disc} for $\theta^n$ and $\omega^n$.
        }
        \titleditem{Update \ac{dom}}{ \label{fl_rigid_step_2}
                Solve  the \acp{bvp} \crefs{eq_domain_fluid_rigid_disc,bc_ela_1_disc,bc_ela_2_disc,eq_domain_fluid_rigid_disc_dot,bc_ela_1_disc_dot,bc_ela_2_disc_dot} for $\budom^n$ and $\budomdot^n$.
        }
        \titleditem{Update \ac{bflu}}{
                \label{fl_rigid_step_3}
                Solve the \acp{bvp} \crefs{eq_fl_rigid_fl_1_reference_disc_aux,bc_fl_rigid_1_reference_disc_aux,bc_fl_rigid_2_reference_disc_aux,bc_fl_rigid_3_reference_disc_aux,bc_fl_rigid_4_reference_disc_aux,bc_fl_rigid_5_reference_disc_aux} and \crefs{eq_phi_2,bc_phi_1,bc_phi_2}  for $\bvbarf$ and $\sigmafr^{n-1/2}$, and obtain $\bm \vfr^n$ from \cref{eq_phi_1}.
        }
\end{enumerate}

The fields ${\bm \vfr}^{n-1}$, $\sigmafr^{n-3/2}$ in the \ac{cur} configuration are then obtained from the respective fields in the \ac{ref} configuration through \cref{eq_vf_ref,eq_sigmaf_ref}.




An example of such dynamics is shown in \cref{fig_fluid_rigid}, where we study the interaction between air (\ac{bflu}) and an aerogel (\ac{rig})---an ultra light solid material involving a large porous structure \cite{kistlerCoherentExpandedAerogels1931}. The channel  dimension are $L = 2.2 \, \met$, $h =  0.41 \, \met$, which have been taken from the \ac{feat2d} DFG 2D-3 benchmark example for a fuid flow around a cylinder \cite{blumFEAT2DFiniteElement1995,schaferBenchmarkComputationsLaminar1996}, see  \cref{fig_ref_curr_rigid}. Given that the materials, such as \ac{bflu} and \ac{bod}, are two dimensional, in what follows we will estimate their parameter values by considering their three-dimensional counterparts, and assume a unit thickness in the out-of-plane direction.
The parameters for \ac{bod} have been estimated from the  properties of air: $\rhof = 1 \, \kg / \met^2$, $\etaf = 10^{-5} \, \pas\, \met \, \sec$ \cite{hilsenrath1955tables}. The  inflow velocity profile is chosen as a Poiseuille flow
\be\label{eq_inflow_profile}
\vfin(\bxc) = \frac{6 v_0}{h^2} \xc^2 (\xc^2 - h),
\ee
with $v_0 = 1 \, \met/\sec$. The \ac{dom} mesh stiffness exponent  has been chosen to be $\zeta = 3$ \cite{kamenskyLectureNotesMAE}.  Finally,  \ac{rig} is an ellipse with semi-axes $a = 0.1\, \met$, $b = 0.075\, \met$,  center located at  $\bxr = (0.25 \, \met, 0.2\, \met)$, and moment of inertia $I =10^{-3} \,  \kg \, \met^2$, which corresponds to an aerogel density of $\sim 6 \, \kg/\met^3$ \cite{niculescuUpdatedOverviewSilica2024}. As shown in \cref{fig_fluid_rigid_timescale},  \ac{rig}  has an oscillatory dynamics, whose period can be understood with scaling arguments---see \cref{sec_fluid_rigid_timescale}.


\begin{figure}
        \centering
        \includegraphics[width=\columnwidth]{figures/figure_8/figure_8.pdf}
        \caption{
                \label{fig_fluid_rigid_timescale}
                Angle of the \acl{rig} as a function of time for the interaction between a  \acl{bflu} and a  \acl{rig} shown in \cref{fig_fluid_rigid}.
        }
\end{figure}



\begin{figure*}
        \centering
        \includegraphics[width=\textwidth]{figures/figure_2/figure_2.pdf}
        \caption{
                \label{fig_fluid_rigid}
                Interaction between a  \acf{bflu} and a rigid \acf{rig}, whose material properties roughly correspond to air (\ac{bflu}) and an aerogel (\ac{rig}).  \plab{A} Temporal snapshot of the mesh (black lines) and  velocity field (arrows) at the \acl{cur} time $t$, shown on top. The direction of the velocity field is indicated by the arrows, and its norm by their color; the direction of arrows with velocity close to zero is not defined. Here,  \ac{rig} (ellipse) is allowed to pivot about its left focal point (red dot), and the related angle with respect to the $\xc^1$ direction is denoted by a red arc.  Here, the Reynolds number is $R \sim 2 \times 10^2$.
                \plab{B} Color map of the surface tension at the \acl{cur} instant of time $t$.
        }
\end{figure*}